
\documentclass[11pt,a4paper]{article}

% ============================================================
% Compile with XeLaTeX
% ============================================================
\usepackage{kotex}

\usepackage{geometry}
\geometry{top=2.8cm,bottom=2.8cm,left=2.5cm,right=2.5cm}

\usepackage{graphicx}
\usepackage{booktabs}
\usepackage{tabularx}
\usepackage{longtable}
\usepackage{array}
\usepackage{multirow}
\usepackage{float}
\usepackage{enumitem}
\usepackage{setspace}
\usepackage{titlesec}
\usepackage{xcolor}
\usepackage{hyperref}
\usepackage{listings}
\usepackage{caption}
\usepackage{amsmath}
\usepackage{amssymb}
\usepackage{url}

\setstretch{1.45}
\setlength{\parindent}{1em}
\setlength{\parskip}{0.25em}

\hypersetup{
    colorlinks=true,
    linkcolor=black,
    citecolor=black,
    urlcolor=blue,
    pdftitle={솔티(Sortie): OCR 및 온디바이스 멀티모달 분류 모델을 활용한 맞춤형 스크린샷 아카이빙 시스템
}
}

\titleformat{\section}{\Large\bfseries}{\thesection.}{0.6em}{}
\titleformat{\subsection}{\large\bfseries}{\thesubsection.}{0.5em}{}
\titleformat{\subsubsection}{\normalsize\bfseries}{\thesubsubsection.}{0.5em}{}

\lstset{
    basicstyle=\ttfamily\small,
    breaklines=true,
    frame=single,
    columns=fullflexible
}

\newcolumntype{Y}{>{\raggedright\arraybackslash}X}
\newcommand{\placeholder}[1]{\fbox{\parbox{0.94\linewidth}{\textbf{[#1]}}}}

\begin{document}

% ============================================================
% 표지
% ============================================================
\begin{titlepage}
\begin{flushright}
\textbf{Pusan National University}\\
\textbf{Department of Computer Science and Engineering}\\
\textbf{Technical Report 2026}
\end{flushright}

\vspace*{3.2cm}

\begin{center}
{\Huge\bfseries 솔티(Sortie)}\\[0.8cm]
{\Large\bfseries OCR 및 온디바이스 멀티모달 분류 모델을 활용한 
}\\
{\Large\bfseries 맞춤형 스크린샷 아카이빙 시스템}

\vspace{2.8cm}

\placeholder{부산대학교 로고 삽입}

\vspace{2.5cm}

{\large
저자 1 학번 성명 메일주소\\[0.3cm]
저자 2 학번 성명 메일주소\\[0.3cm]
저자 3 학번 성명 메일주소
}

\vspace{1.4cm}
{\large 지도교수 성명}
\end{center}

\vfill
\begin{center}
2026년 9월
\end{center}
\end{titlepage}

\tableofcontents
\newpage

% ============================================================
% 초록
% ============================================================
\section*{초록}
\addcontentsline{toc}{section}{초록}

스마트폰 사용 과정에서 생성되는 스크린샷은 일정, 장소, 상품, 메모와 같이 서로 다른 종류의 정보를 포함하지만, 일반적인 갤러리 환경에서는 대부분 이미지 파일로만 저장된다. 사용자가 이러한 정보를 다시 활용하기 위해서는 이미지를 찾아 열고, 필요한 내용을 직접 확인한 뒤 캘린더나 메모 애플리케이션 등에 다시 입력해야 한다. 본 연구에서는 이러한 비정형 스크린샷을 사용자가 활용할 수 있는 구조화된 정보로 변환하는 모바일 애플리케이션 ``솔티(Sortie)''를 구현하였다.

솔티는 Flutter 기반 모바일 클라이언트와 Python 기반 백엔드 서버를 결합한 하이브리드 구조를 사용한다. 클라이언트에서는 Google ML Kit Text Recognition을 이용하여 이미지에서 OCR 텍스트를 추출하고, TFLite 기반 온디바이스 텍스트 분류기를 이용하여 SCHEDULE, PLACE, WISHLIST, MEMO의 네 가지 대분류를 판단한다. 분류 모델의 입력은 프로젝트에 포함된 vocabulary와 IDF 값을 이용해 구성한 TF-IDF 벡터이며, 단어 경계 기반 character bigram 및 trigram을 특징으로 사용한다. 이후 개인정보 보호를 위해 온디바이스 NER 모델과 정규표현식 기반 탐지 로직을 적용하고, 탐지된 민감정보는 서버로 전송되기 전에 AES-256-CBC 방식으로 암호화하여 \texttt{[ENC:...]} 형태의 토큰으로 치환한다.

서버에서는 FastAPI를 기반으로 분석 요청을 처리하며, 카테고리별 프롬프트를 사용하여 클라우드 LLM에 세부 정보 추출을 요청한다. LLM 응답은 검증 모듈을 통해 타입, 필수 필드, confidence 및 카테고리별 부가 필드가 확인되고 보강된다. 또한 다중 요청에 대한 동시 LLM 호출 수를 3개로 제한하는 Semaphore, 이미지 해시 기반 중복 결과 조회, 비동기 일괄 처리, 알림 스케줄러, iCalendar 생성 및 카카오톡 캘린더 연동 등 후속 활용을 위한 기능을 구현하였다.

실험에서는 프로젝트에 포함된 \texttt{evaluate\_extraction.py}를 이용하여 실제 사용 데이터 673건 중 카테고리별 50건씩 총 200건을 샘플링하고, GPT-4o가 생성한 기준값과 온디바이스 Regex 방식 및 GPT-4o-mini 방식의 정보 추출 결과를 비교하였다. 전체 평균 측정 시간은 Regex 0.135 ms, LLM 1.06 s였으며, 장소명이나 상품명과 같은 비정형 필드에서는 LLM이 Regex보다 높은 매칭률을 보였다. 반면 날짜나 가격처럼 정형적인 패턴은 Regex가 경쟁력 있는 결과를 보였다. 이러한 결과는 정형 정보의 빠른 로컬 처리와 문맥 기반 정보의 LLM 보강을 결합하는 현재의 하이브리드 구조가 서로 다른 정보 유형의 특성을 고려한 설계임을 보여준다. 다만 해당 실험의 기준값은 사람의 독립적인 정답 라벨이 아니라 GPT-4o의 출력이며, 본 연구에서는 이를 절대적 ground truth가 아닌 ``reference annotation''으로 해석한다.

\textbf{주제어:} 스크린샷 구조화, OCR, 온디바이스 AI, TF-IDF, Character N-gram, NER, 개인정보 보호, LLM, FastAPI, Flutter

\newpage

% ============================================================
% 1 서론
% ============================================================
\section{서론}

\subsection{연구 배경}

스마트폰은 일상적인 정보 획득과 기록을 위한 핵심 도구가 되었다. 특히 웹페이지, 메신저, 지도, 쇼핑몰, 예약 서비스 등에서 확인한 정보를 빠르게 보관하기 위해 스크린샷 기능이 널리 사용된다. 사용자는 ``나중에 갈 장소'', ``구매하고 싶은 상품'', ``예약 일정'', ``기억해야 할 텍스트'' 등을 별도의 기록 절차 없이 이미지로 남길 수 있다.

그러나 스크린샷은 기록이 쉽다는 장점과 동시에 정보 재활용이 어렵다는 문제를 가진다. 이미지 내부에 날짜와 시간, 장소명, 상품명, 가격, 주소 등의 정보가 포함되어 있어도 파일 시스템 수준에서는 하나의 이미지로만 취급되기 때문이다. 시간이 지나면 사용자는 수많은 이미지 중 필요한 스크린샷을 다시 찾아야 하며, 발견한 정보를 다시 읽고 필요한 애플리케이션에 옮겨 적어야 한다.

기존 OCR 기술을 이용하면 이러한 이미지에서 텍스트를 추출할 수 있지만, 단순 OCR은 ``문자열을 얻는 것''과 ``정보를 이해하여 구조화하는 것'' 사이의 차이를 해결하지 못한다. 예를 들어 ``체크인 8월 20일 15:00 / 체크아웃 8월 22일 11:00''이라는 텍스트에서 사용자가 실제로 필요로 하는 정보는 단순한 문자 배열이 아니라 하나의 예약 일정에 연결된 시작일, 종료일, 시작 시간, 종료 시간이다.

본 연구는 이 문제를 스크린샷을 ``이미지''가 아니라 ``잠재적인 구조화 데이터''로 보는 관점에서 해결하고자 한다. 솔티는 OCR을 통해 텍스트를 확보하고, 온디바이스 모델과 규칙 기반 처리를 이용하여 빠른 전처리를 수행한 후, 필요한 의미적 해석은 서버의 LLM에 위임한다. 또한 클라우드 분석 이전에 개인정보를 보호하기 위한 별도의 전처리 계층을 배치하였다.

\subsection{기존 문제점}

본 프로젝트가 해결하고자 한 문제는 크게 네 가지로 정리할 수 있다.

첫째, 스크린샷 정보의 검색성과 재사용성이 낮다. 이미지 갤러리에서는 스크린샷에 포함된 의미 정보를 기준으로 검색하거나 분류하기 어렵다.

둘째, OCR만으로는 구조화가 충분하지 않다. OCR은 이미지 속 문자를 추출하지만 날짜, 장소, 상품, 메모 등의 의미적 관계를 자동으로 구성하지 않는다.

셋째, 클라우드 LLM을 직접 사용하는 경우 개인정보가 포함된 원문을 서버로 전송해야 할 가능성이 있다. 실제 스크린샷에는 전화번호, 카드번호, 계좌번호, 주민등록번호, 여권번호, 예약번호와 같은 민감정보가 포함될 수 있다.

넷째, 모든 작업을 서버에서 수행할 경우 네트워크 지연과 API 호출 비용이 발생한다. 반대로 모든 작업을 온디바이스에서 수행하면 자연어 문맥을 이해해야 하는 장소명, 상품명, 요약 등의 처리에서 한계가 발생할 수 있다.

따라서 본 프로젝트에서는 ``모든 것을 로컬에서 처리''하거나 ``모든 것을 서버에서 처리''하는 단일 방식보다 각 처리의 특성에 맞추어 로컬과 서버의 역할을 분리하는 방식을 선택하였다.

\subsection{연구 목표}

본 연구의 목표는 다음과 같다.

\begin{enumerate}[label=\arabic*)]
    \item 스크린샷에서 OCR을 통해 텍스트를 자동 획득한다.
    \item 획득한 텍스트를 네 가지 대분류로 분류하여 후속 처리를 단순화한다.
    \item 서버로 전송하기 전에 온디바이스에서 민감정보를 탐지하고 암호화한다.
    \item 카테고리별로 필요한 필드를 LLM을 이용해 구조화한다.
    \item 구조화된 결과를 로컬 DB에 저장하고 검색, 수정, 확인할 수 있도록 한다.
    \item 일정, 알림, 캘린더 등 구조화된 데이터의 후속 활용을 지원한다.
    \item 실제 코드와 평가 데이터에 근거하여 시스템의 성능과 한계를 분석한다.
\end{enumerate}

\subsection{연구 범위 및 차별점}

본 연구는 일반적인 문서 OCR 시스템을 만드는 것이 아니라, 스마트폰에서 발생하는 ``개인적인 스크린샷''을 대상으로 한다. 따라서 시스템의 핵심은 OCR 정확도 자체보다 OCR 이후의 의미 분류, 구조화, 개인정보 보호, 저장 및 재활용을 하나의 파이프라인으로 연결하는 데 있다.

또한 프로젝트는 개인정보를 단순히 별표 형태로 삭제하는 방식이 아니라, 온디바이스에서 식별한 민감정보를 암호화된 토큰으로 변환하여 후속 LLM 분석에서 문장의 나머지 구조를 유지하도록 설계하였다. 이는 민감정보가 존재하더라도 서버 분석 자체를 포기하지 않고, 필요한 문맥을 최대한 유지하려는 설계이다.

\newpage

% ============================================================
% 2 연구 배경
% ============================================================
\section{연구 배경}

\subsection{스크린샷 기반 비정형 데이터}

스크린샷은 생성 과정에서 별도의 데이터 스키마를 갖지 않는다. 같은 이미지 안에 날짜, 가격, 장소, 설명, 버튼, 메뉴, 광고, 상태바 등이 동시에 존재할 수 있다. 특히 모바일 화면은 본문과 UI 요소가 섞여 있어 OCR 결과에도 불필요한 문자열이 포함될 가능성이 있다.

따라서 스크린샷 처리에서는 다음 세 단계가 필요하다.

\begin{enumerate}
    \item 이미지에서 텍스트를 추출하는 OCR 단계
    \item 텍스트의 의미적 종류를 판단하는 분류 단계
    \item 카테고리별 필드로 변환하는 구조화 단계
\end{enumerate}

솔티는 이 세 단계에 개인정보 보호 단계를 추가하여 ``OCR $\rightarrow$ 분류 $\rightarrow$ 민감정보 보호 $\rightarrow$ 구조화''의 파이프라인을 구성한다.

\subsection{OCR 기술 선택}

프로젝트의 Flutter 클라이언트는 \texttt{google\_mlkit\_text\_recognition} 패키지를 사용한다. OCR을 모바일 클라이언트에서 수행함으로써 이미지 전체를 서버에 전송하지 않고 텍스트 추출을 먼저 수행할 수 있다. 이 구조는 네트워크 전송량을 줄이는 동시에 개인정보 보호 계층을 서버 요청보다 앞에 배치할 수 있다는 장점이 있다.

OCR 결과는 단순 문자열이 아니라 이후 분류 및 마스킹의 입력으로 사용된다. 실제 코드에서는 상단 상태바의 시간, 날짜, 통신사 등의 문자열이 본문에 혼입되는 것을 줄이기 위해 OCR 텍스트의 초기 두 줄을 대상으로 상태바 패턴을 검사하여 제거하는 로직도 포함한다.

\subsection{온디바이스 텍스트 분류}

프로젝트의 \texttt{OnDeviceTextClassifier}는 TFLite Interpreter를 이용하여 \texttt{classifier\_demo.tflite}를 로드한다. 함께 제공되는 \texttt{vocab.json}, \texttt{idf.json}, \texttt{classifier\_label\_map.json}을 이용하여 텍스트를 모델 입력으로 변환한다.

전처리는 다음과 같이 구현되어 있다.

\begin{enumerate}
    \item 한글, 영문, 숫자를 제외한 특수문자를 공백으로 치환
    \item 연속된 공백을 하나로 정규화
    \item 소문자 변환
    \item 공백을 기준으로 단어 분리
    \item 각 단어에 공백을 붙인 뒤 character bigram과 trigram 생성
    \item 프로젝트 vocabulary에 존재하는 n-gram만 TF-IDF 벡터에 반영
    \item L2 normalization 수행
    \item TFLite 모델에 입력하여 네 개 클래스의 확률 계산
    \item 최댓값의 클래스를 최종 카테고리로 선택
\end{enumerate}

이 방식은 형태소 분석기와 같은 별도의 무거운 언어 처리기를 사용하지 않고, 모델 학습 시 사용한 특징 추출 방식을 모바일 코드에 이식한다는 특징이 있다. Character N-gram은 단어 전체가 정확히 일치하지 않아도 부분적인 문자열 특징을 사용할 수 있으므로 OCR 오류나 형태 변화에 대한 보완 수단으로 활용될 수 있다.

\subsection{TF-IDF와 Character N-gram}

TF-IDF는 문서 내에서 자주 나타나지만 전체 문서 집합에서는 상대적으로 희소한 단어 또는 특징에 더 높은 가중치를 부여하는 방식이다. 본 프로젝트에서는 일반적인 단어 단위 TF-IDF가 아니라 단어 경계를 고려한 character bigram과 trigram을 사용한다.

문서 $d$에서 특징 $t$의 TF-IDF는 개념적으로 다음과 같이 표현할 수 있다.

\[
\mathrm{TFIDF}(t,d)=\mathrm{TF}(t,d)\times \mathrm{IDF}(t)
\]

실제 모바일 코드에서는 프로젝트의 \texttt{idf.json}에 저장된 IDF 값을 불러와 각 n-gram의 등장 횟수와 곱한다. 이후 전체 벡터를 L2 정규화한다.

이러한 설계는 모델과 모바일 전처리의 불일치를 줄인다는 점에서 중요하다. 서버에서 사용하는 별도의 자연어 처리 방식이 아니라 학습 시 사용한 특징을 동일한 형태로 생성하기 때문에 TFLite 모델이 기대하는 입력 공간을 유지할 수 있다.

\subsection{NER와 개인정보 보호}

정규표현식은 일정한 패턴을 가진 개인정보 탐지에 효과적이지만, 사람 이름과 같이 문맥에 의존하는 정보에는 한계가 있다. 반대로 NER 모델은 문맥을 고려할 수 있지만 모든 번호 형식에 대해 동일한 강점을 가지지는 않는다.

솔티의 \texttt{MaskingHelper}는 이 두 방식을 결합한다. 먼저 온디바이스 NER 모델이 WordPiece 기반 토크나이저를 통해 최대 128개 토큰 길이의 입력을 구성하고 TFLite 추론을 수행한다. 현재 코드에서 마스킹 대상으로 사용하는 NER 엔티티는 \texttt{PER}와 \texttt{PHONE}이다.

이후 Regex 계층에서는 주민등록번호, 카드번호, CVC/CVV, 쿠폰 번호, 계좌번호, 전화번호, 예약번호, 여권번호 및 MRZ, 운전면허번호, 신분증 주소와 날짜 등 여러 패턴을 별도로 처리한다. 즉, ``AI가 전부 탐지한다''는 구조가 아니라 NER와 패턴 기반 탐지를 역할별로 결합한 구조이다.

\subsection{암호화 방식}

민감정보의 치환에는 \texttt{CryptoHelper}가 사용된다. 코드에서는 32바이트의 랜덤 키를 생성하여 \texttt{FlutterSecureStorage}에 저장하고, AES-256-CBC 방식으로 문자열을 암호화한다. 각 암호화마다 16바이트의 랜덤 IV를 생성하고 IV와 암호문을 결합한 뒤 Base64로 인코딩한다. 결과는 다음과 같은 형태의 문자열로 반환된다.

\begin{lstlisting}
[ENC:Base64EncodedIVAndCiphertext]
\end{lstlisting}

서버에 전달되는 OCR 텍스트에는 이러한 토큰이 포함될 수 있으며, 클라이언트의 복호화 함수는 \texttt{[ENC:...]} 패턴을 찾아 로컬에서 다시 평문으로 복원한다.

다만 AES-CBC 자체는 암호화만 제공하며 메시지 인증을 내장하지 않는다. 따라서 본 구현을 ``현대적인 의미의 완전한 authenticated encryption''으로 확대 해석해서는 안 된다. 본 프로젝트의 구현은 민감 문자열을 서버 요청에서 평문으로 보내지 않기 위한 보호 계층으로 기술하는 것이 정확하며, 향후에는 AES-GCM과 같은 인증 암호화 방식으로 개선할 수 있다.

\subsection{LLM과 Structured Output}

스크린샷 정보 중 날짜나 가격처럼 정형적인 정보는 정규표현식으로 처리하기 쉽지만, ``이 장소가 어느 식당인지'', ``이 화면에서 상품명이 무엇인지'', ``이 메모의 핵심 제목이 무엇인지''와 같은 정보는 문맥 이해가 필요하다.

따라서 서버에서는 카테고리별 프롬프트를 사용하여 LLM에게 필요한 필드를 추출하도록 한다. SCHEDULE, PLACE, WISHLIST, MEMO마다 서로 다른 필드 정의와 예외 규칙을 사용한다.

예를 들어 SCHEDULE 프롬프트에서는 체크인과 체크아웃을 하나의 일정으로 연결하고, 기프티콘의 유효기간을 \texttt{expires\_at}에 저장하며, 취소 마감일을 일반 일정의 시작일이나 만료일로 오인하지 않도록 명시한다. PLACE에서는 여러 장소가 나열될 경우 각각을 객체로 분리하고, WISHLIST에서는 여러 상품을 별도의 객체로 분리하도록 지시한다.

\subsection{SQLite와 JSON}

모바일 클라이언트와 서버 모두 SQLite를 사용한다. 클라이언트의 \texttt{screenshots} 테이블은 \texttt{id}, \texttt{asset\_id}, \texttt{type}, \texttt{confidence}, \texttt{fields}, \texttt{image\_path}, \texttt{status}, \texttt{created\_at}, \texttt{updated\_at} 등의 정보를 저장한다.

세부 추출 필드는 카테고리에 따라 달라지므로 고정 컬럼 대신 \texttt{fields}에 JSON 문자열을 저장한다. 서버 DB에는 이미지 해시를 저장하여 동일 이미지의 분석 결과를 조회할 수 있도록 하며, 알림 기능을 위해 기기 토큰과 알림 내역 테이블도 존재한다.

\subsection{선행연구 및 기술 선택의 시사점}

기존 초안에서 참고한 선행연구는 개인정보 보호형 LLM 처리, 온디바이스 모델의 한계, 프롬프트 엔지니어링 등의 관점에서 본 프로젝트의 설계 방향을 설명하는 데 활용할 수 있다. 다만 최종보고서에서는 실제로 확인되지 않은 성능이나 특정 연구의 결론을 프로젝트 결과처럼 서술하지 않고, ``기술 선택에 참고한 배경''으로 한정한다.

\begin{table}[H]
\centering
\small
\begin{tabularx}{\linewidth}{p{3.0cm}Y Y}
\toprule
\textbf{기술} & \textbf{프로젝트에서의 역할} & \textbf{선택 이유}\\
\midrule
ML Kit OCR & 이미지 $\rightarrow$ 텍스트 & 모바일에서 직접 수행하여 서버 전송 전 전처리 가능\\
TFLite 분류기 & 네 가지 대분류 & 모바일에서 빠른 추론 및 네트워크 비의존적 분류 가능\\
Character N-gram + TF-IDF & 분류 입력 특징 & 프로젝트의 학습 특징을 모바일에서 재현\\
NER + Regex & 민감정보 탐지 & 문맥 기반 탐지와 명시적 패턴 탐지를 결합\\
AES-256-CBC & 민감 문자열 보호 & 서버 전송 전 원문 값을 암호문 토큰으로 치환\\
LLM & 상세 필드 추출 & 비정형 정보의 문맥 기반 구조화\\
SQLite + JSON & 결과 저장 & 카테고리별 상이한 필드를 유연하게 저장\\
FastAPI & 분석 API & 모바일 클라이언트와 LLM 처리 계층 분리\\
\bottomrule
\end{tabularx}
\caption{주요 기술의 선택 및 역할}
\end{table}

\newpage

% ============================================================
% 3 연구 내용
% ============================================================
\section{연구 내용}

\subsection{전체 시스템 아키텍처}

솔티의 전체 처리 흐름은 모바일 클라이언트와 서버의 역할을 분리한 구조로 구성된다.

\placeholder{그림 삽입: 전체 시스템 아키텍처 -- 갤러리 이미지 $\rightarrow$ ML Kit OCR $\rightarrow$ 온디바이스 분류/마스킹 $\rightarrow$ FastAPI $\rightarrow$ LLM $\rightarrow$ Validator $\rightarrow$ SQLite/앱 UI}

사용자가 이미지를 선택하면 Flutter 클라이언트가 이미지에서 OCR 텍스트를 획득한다. 이후 로컬 분류기가 텍스트를 SCHEDULE, PLACE, WISHLIST, MEMO 중 하나로 분류한다. 개인정보 보호가 필요한 텍스트는 NER 및 Regex 계층을 통과하며 민감한 구간은 암호화 토큰으로 변환된다.

그 다음 클라이언트는 분류 결과와 마스킹된 OCR 텍스트를 FastAPI 서버의 분석 API에 전달한다. 서버는 카테고리별 프롬프트와 입력 텍스트를 LLM에 전달하고, 반환된 JSON을 검증 및 보강한 뒤 클라이언트에 반환한다.

\subsection{모바일 클라이언트}

\subsubsection{이미지 수집과 OCR}

프로젝트는 \texttt{image\_picker}, \texttt{photo\_manager}, \texttt{wechat\_assets\_picker}, \texttt{image\_cropper} 등의 패키지를 사용한다. 이를 통해 갤러리 이미지 선택과 이미지 관련 작업을 수행한다.

OCR은 \texttt{google\_mlkit\_text\_recognition}을 이용한다. OCR 결과는 화면에서 초안 데이터로 관리되며 이후 온디바이스 분류 및 AI 분석 요청의 입력으로 사용된다.

상단 상태바 제거 로직은 OCR 결과에 포함될 수 있는 현재 시간, 날짜, 통신사 등의 정보를 본문 정보와 구분하기 위한 것이다. 이 단계는 정보 추출 정확도를 높이기 위한 단순한 노이즈 제거 단계로 볼 수 있다.

\subsubsection{온디바이스 대분류}

온디바이스 분류기는 다음 네 가지 클래스를 반환한다.

\begin{table}[H]
\centering
\begin{tabularx}{0.8\linewidth}{c l Y}
\toprule
\textbf{Index} & \textbf{Type} & \textbf{의미}\\
\midrule
0 & SCHEDULE & 일정, 예약, 기한 등 시간 중심 정보\\
1 & PLACE & 장소, 맛집, 위치 등 장소 중심 정보\\
2 & WISHLIST & 상품, 쇼핑, 구매 관심 정보\\
3 & MEMO & 기타 기록 및 참고 정보\\
\bottomrule
\end{tabularx}
\caption{온디바이스 분류기의 대분류}
\end{table}

모델이 초기화되지 않았거나 전처리 결과가 비어 있는 경우 코드의 폴백 값은 MEMO이다. 또한 주민등록증, 운전면허증, 여권 등의 신분증 텍스트가 감지되는 경우에는 민감한 문서의 내용을 일반적인 일정이나 장소로 잘못 분류하지 않도록 MEMO 카테고리로 강제하는 로직이 존재한다.

\subsubsection{기프티콘 사전 처리}

프로젝트에는 기프티콘 관련 OCR을 처리하기 위한 별도의 휴리스틱이 존재한다. 기프티콘, 쿠폰, 바코드, 교환권, 모바일상품권, 선물하기, 교환처 등의 키워드를 이용하여 해당 문서의 특성을 판단한다.

코드에는 유효기간, 교환처, 상품명을 탐색하는 \texttt{\_normalizeGifticon} 함수가 존재하지만, 실제 최종 마스킹 흐름에서는 원본 OCR 텍스트를 3줄로 임의 축약하면 상품명이나 상호명이 손실될 수 있다는 문제를 확인하여 해당 축약 결과를 LLM 입력으로 사용하지 않도록 한 수정 내용도 포함되어 있다. 이 부분은 개발 과정에서 발견한 문제를 코드 수준에서 수정한 사례로 기술할 수 있다.

\subsection{개인정보 보호 파이프라인}

\subsubsection{1단계: NER 기반 탐지}

\texttt{NerClassifier}는 앱 시작 시 \texttt{ner\_model.tflite}, \texttt{vocab.txt}, \texttt{ner\_label\_map.json}을 로드한다. 토크나이저는 WordPiece 방식으로 단어를 분해하고 최대 128개 토큰 길이에 맞춘다.

모델은 각 토큰에 대해 라벨을 출력하며, 현재 코드에서 \texttt{B-PER}, \texttt{I-PER}, \texttt{B-PHONE}, \texttt{I-PHONE} 등의 라벨 중 PER과 PHONE 엔티티가 마스킹 대상이다. 탐지된 토큰의 원문 offset을 이용하여 해당 문자열 범위를 찾고, 겹치거나 인접한 범위를 병합한 후 역순으로 암호화한다.

\subsubsection{2단계: Regex 기반 탐지}

NER 이후에는 여러 정형 패턴을 Regex로 검사한다.

\begin{longtable}{p{3.4cm}p{7.8cm}}
\toprule
\textbf{대상} & \textbf{구현 특징}\\
\midrule
주민등록번호 & 6자리 생년월일 + 1--8로 시작하는 7자리 패턴\\
카드번호 & 공백/하이픈을 고려한 14--16자리 계열 패턴 및 카드 키워드 검증\\
CVC/CVV & CVC, CVV, 보안코드 등의 키워드와 3--4자리 숫자 결합\\
쿠폰번호 & 12--16자리 숫자와 기프티콘 관련 키워드의 결합\\
계좌번호 & 여러 구간의 숫자와 하이픈 패턴, 은행 관련 키워드 검증\\
전화번호 & 01X 및 지역번호 기반 전화번호 패턴\\
예약번호 & 예매번호, 예약번호, 티켓번호 등의 키워드와 영숫자 조합\\
여권번호 & 여권 키워드 및 영숫자 번호 패턴\\
여권 MRZ & ``<'' 문자를 포함하는 판독 영역 후보를 인접 줄과 함께 탐지\\
운전면허번호 & 지역 코드 및 숫자 조합 패턴\\
신분증 주소/날짜 & 신분증 문서가 확인된 경우 주소 및 날짜 패턴을 추가 처리\\
\bottomrule
\end{longtable}

\subsubsection{암호화 및 서버 전달}

탐지된 값은 \texttt{CryptoHelper().encryptSensitive()}를 통해 암호화된다. 따라서 LLM 요청의 텍스트에는 원래의 번호 대신 \texttt{[ENC:...]} 형태의 토큰이 들어간다.

클라이언트가 서버로 전달하는 분석 요청에는 현재 날짜와 시스템 지시문도 포함된다. 특히 \texttt{[ENC:...]} 토큰을 임의로 해석하거나 변경하지 말고 그대로 JSON 값으로 유지하라는 규칙이 프롬프트에 포함되어 있다.

이 구조의 중요한 점은 ``LLM이 개인정보를 삭제한다''가 아니라, \textbf{LLM에 도달하기 전에 클라이언트가 민감 문자열을 변환한다}는 것이다.

\subsection{서버 시스템}

\subsubsection{FastAPI 분석 API}

백엔드에는 FastAPI 애플리케이션이 구성되어 있으며 분석 관련 라우터와 알림, 결과 조회 등의 기능이 분리되어 있다. 모바일 클라이언트는 \texttt{/api/analyze/v2} 엔드포인트를 이용하여 OCR 텍스트와 분류 결과를 전달한다.

단건 요청은 JSON 형태로 다음과 같은 개념적 구조를 갖는다.

\begin{lstlisting}
{
  "ocr_text": "... masked OCR text ...",
  "type": "SCHEDULE",
  "masked_tokens": []
}
\end{lstlisting}

다중 요청에서는 \texttt{items} 배열을 사용하여 여러 분석 항목을 한 번에 전달한다.

\subsubsection{LLM 호출과 재시도}

서버의 LLM 클라이언트는 모델 호출을 별도 모듈로 분리한다. LLM 호출은 최대 재시도 횟수를 고려하는 구조를 갖고 있으며, 호출이 최종적으로 실패하면 MEMO, confidence 0, 빈 fields 및 수동 입력 필요 메시지를 포함한 폴백 결과를 반환한다.

단, 보고서에서는 실제 API 공급자나 모델 버전이 변경될 수 있으므로 ``현재 평가에서는 GPT-4o-mini를 사용했다''와 ``시스템은 LLM 호출 계층을 추상화한다''를 구분하여 서술한다.

\subsubsection{동시성 제어}

서버에는 \texttt{MAX\_CONCURRENT\_LLM = 3}으로 정의된 Semaphore가 있다. 여러 요청이 동시에 들어오더라도 동시에 실행되는 LLM 호출 수를 제한하여 API Rate Limit이나 서버 측 자원 과부하 가능성을 줄이는 역할을 한다.

\[
C_{\max}=3
\]

여기서 $C_{\max}$는 동시에 LLM API를 호출할 수 있는 최대 수이다. 실제 처리 시간은 네트워크 상태와 LLM 응답 시간에 따라 달라지므로 이 값 자체를 ``20건 처리 시간이 항상 일정하다''는 의미로 해석해서는 안 된다.

\subsection{카테고리별 정보 구조화}

\subsubsection{SCHEDULE}

SCHEDULE은 시간에 따라 행동이 필요한 정보를 표현한다. 서버 프롬프트와 validator는 다음과 같은 필드를 다룬다.

\begin{itemize}
    \item \texttt{title}
    \item \texttt{start\_at}, \texttt{end\_at}
    \item \texttt{expires\_at}
    \item \texttt{start\_time}, \texttt{end\_time}
    \item \texttt{exchange\_place}
    \item \texttt{participants}
    \item \texttt{recurrence}
    \item \texttt{cancellation\_deadline}
    \item \texttt{sub\_type}
\end{itemize}

세부 유형은 GIFTICON, SUBSCRIPTION, APPOINTMENT, TICKET, DEADLINE, DELIVERY 등을 사용한다. validator는 유형에 따라 \texttt{keep\_photo}, \texttt{reminder\_days}, \texttt{calendar\_type} 등의 값을 보강한다.

\subsubsection{PLACE}

PLACE는 장소의 이름과 지역 정보를 중심으로 구조화된다. 서버의 후처리에서는 지역명을 15개 표준 광역 단위로 정규화하고, ``해운대''가 포함된 경우 부산으로 처리하는 예외도 존재한다.

또한 \texttt{map\_ready}는 주소 또는 지역 정보가 존재하는지를 이용하여 지도 관련 UI를 사용할 수 있는지 판단한다. 코드에서는 초기 프롬프트에 존재했던 평점 필드를 최종 enrichment 과정에서 제거한다. 따라서 최종 구현을 기술할 때는 ``평점을 저장한다''고 서술하지 않는다.

\subsubsection{WISHLIST}

WISHLIST는 상품 단위의 정보로 구조화된다. 주요 필드는 상품명, 가격, 옵션, 판매처, URL이다. 여러 상품이 하나의 OCR 텍스트에 포함될 경우 서버 프롬프트는 각각의 객체로 분리하도록 설계되어 있다.

validator의 후처리에서는 가격 문자열에서 숫자와 소수점을 제외한 문자를 제거하여 수치형으로 정규화한다. \texttt{keep\_photo}의 기본값은 true이다.

\subsubsection{MEMO}

MEMO는 일정, 장소, 상품에 해당하지 않는 일반적인 참고 정보를 대상으로 한다. 서버 프롬프트에서는 RECIPE, NOVEL, CHECKLIST, ARTICLE, NOTE, OTHER 등의 \texttt{sub\_type}을 사용한다.

예를 들어 RECIPE에서는 재료와 조리 순서, NOVEL에서는 작품명과 작가, CHECKLIST에서는 체크 항목 등의 필드를 추가로 다룬다. title이 존재하지 않는 경우 body 앞부분을 이용해 제목을 생성하는 후처리도 구현되어 있다.

\subsection{Validator와 데이터 품질 관리}

LLM은 확률적인 생성 모델이므로 JSON 형식이 맞더라도 필수 정보가 누락될 수 있다. 이를 보완하기 위해 validator가 별도로 존재한다.

검증 단계에서는 다음을 수행한다.

\begin{enumerate}
    \item type이 허용된 네 가지 값인지 검사
    \item confidence가 0과 1 사이인지 검사
    \item 카테고리별 필수 필드 존재 여부 검사
    \item 여러 객체가 반환되는 경우 각각 검사
    \item SCHEDULE의 시작일/만료일 존재 여부 검사
    \item PLACE의 name 또는 region 존재 여부 검사
    \item 카테고리별 enrichment 수행
    \item 누락된 필드를 \texttt{missing\_fields}에 기록
\end{enumerate}

따라서 서버의 역할은 단순히 LLM 결과를 그대로 반환하는 것이 아니라 ``LLM 생성 $\rightarrow$ 구조 검증 $\rightarrow$ 후처리 및 보강 $\rightarrow$ API 응답''으로 확장된다.

\subsection{로컬 DB와 저장 구조}

모바일 DB의 핵심 테이블은 \texttt{screenshots}이다.

\begin{table}[H]
\centering
\small
\begin{tabularx}{\linewidth}{l l Y}
\toprule
\textbf{필드} & \textbf{타입} & \textbf{역할}\\
\midrule
id & INTEGER & 자동 증가 기본키\\
asset\_id & TEXT & 갤러리 원본 식별자\\
type & TEXT & SCHEDULE/PLACE/WISHLIST/MEMO\\
confidence & REAL & 분석 신뢰도 저장\\
fields & TEXT & 카테고리별 JSON 데이터\\
image\_path & TEXT & 보존된 이미지 경로\\
status & TEXT & DRAFT 등의 상태 관리\\
created\_at & TEXT & 생성 시각\\
updated\_at & TEXT & 수정 시각\\
\bottomrule
\end{tabularx}
\caption{모바일 screenshots 테이블 구조}
\end{table}

서버 DB에는 별도의 \texttt{image\_hash} 컬럼이 존재하며, 동일 이미지에 대한 분석 결과를 조회하는 함수가 구현되어 있다. 이를 통해 분석 결과 캐시 및 중복 분석 방지에 활용할 수 있는 기반을 제공한다.

\subsection{검색 및 필터링}

서버의 결과 조회 기능은 type, status, region, category 및 자유 텍스트 q를 기준으로 결과를 검색할 수 있다. JSON으로 저장된 \texttt{fields} 내부 값에 대해 SQLite의 \texttt{json\_extract}와 LIKE 검색을 사용하는 방식이 구현되어 있다.

이는 단순히 ``DB에 결과를 저장한다''는 수준을 넘어 구조화 결과가 실제 애플리케이션의 검색 기준으로 사용될 수 있도록 한 것이다.

\subsection{이미지 저장 및 생명주기}

앱 내부에는 SCHEDULE, PLACE, WISHLIST, MEMO별 디렉터리를 생성하는 \texttt{AppStorage}가 존재한다. 서버의 validator는 카테고리와 SCHEDULE의 세부 유형에 따라 \texttt{keep\_photo} 기본값을 결정한다.

\begin{table}[H]
\centering
\small
\begin{tabularx}{0.9\linewidth}{l l l}
\toprule
\textbf{유형} & \textbf{keep\_photo 기본값} & \textbf{의도}\\
\midrule
GIFTICON & true & 바코드/쿠폰 이미지 보존\\
TICKET & true & 예매 관련 원본 보존\\
WISHLIST & true & 상품 이미지 보존\\
PLACE & false & 텍스트 구조화 후 이미지 의존도 감소\\
MEMO & false & 텍스트 중심 정보\\
APPOINTMENT & false & 일정 구조화 중심\\
SUBSCRIPTION & false & 갱신 정보 구조화 중심\\
DEADLINE & false & 마감 정보 구조화 중심\\
DELIVERY & false & 배송 일정 구조화 중심\\
\bottomrule
\end{tabularx}
\caption{코드에 정의된 keep\_photo 기본 정책}
\end{table}

\textbf{주의:} 위 표는 validator에 정의된 기본 정책을 나타내며, ``모든 원본 파일이 자동으로 물리 삭제된다''는 의미는 아니다. 최종 보고서에서는 실제 이미지 삭제 호출 경로를 별도로 검증한 후 자동 삭제 범위를 확정해야 한다.

\subsection{캘린더 및 알림}

서버에는 iCalendar 형식의 \texttt{.ics} 문자열을 생성하는 기능이 구현되어 있다. SCHEDULE의 날짜와 시간을 ICS 형식으로 변환하여 VEVENT를 구성하고, 제목과 설명 및 위치 정보를 포함할 수 있다.

모바일 클라이언트에는 Kakao Calendar 서비스도 구현되어 있다. Kakao 로그인 및 \texttt{talk\_calendar} 권한을 확인한 후 일정 생성 API를 호출하는 방식이다.

또한 서버에는 알림 스케줄러가 존재한다. 스케줄러는 60초 간격으로 DB를 확인하여 오늘 또는 내일이 대상인 SCHEDULE/GIFTICON 항목을 찾아 알림 내역을 생성한다. 중복 알림 여부도 DB에서 확인한다.

\subsection{다중 이미지 처리}

모바일 클라이언트는 여러 스크린샷을 선택하여 일괄 분석할 수 있으며, 코드에서는 요청을 20개 단위로 분할한다.

\begin{lstlisting}
for (int i = 0; i < selected.length; i += 20) {
    final chunk = selected.sublist(
        i,
        (i + 20 > selected.length)
            ? selected.length
            : i + 20
    );
    ...
}
\end{lstlisting}

서버 측에서는 동시에 실행되는 LLM 호출 수를 3개로 제한한다. 따라서 ``20장의 이미지를 한 번에 LLM에 동시에 무제한 전송''하는 방식이 아니라, 클라이언트의 배치 크기와 서버의 동시성 제한을 함께 사용하는 구조이다.

\subsection{사용자 검토 및 상태 관리}

DB의 \texttt{status} 필드는 DRAFT 등의 상태를 표현하며, 모바일 DB에는 사용자가 확인한 결과를 CONFIRMED 상태로 저장하는 흐름이 구현되어 있다. 또한 \texttt{updateFields}, \texttt{updateTypeAndFields}, \texttt{updateStatus} 등의 CRUD 함수가 존재하여 사용자가 구조화 결과를 수정하거나 상태를 변경할 수 있다.

이 구조는 완전 자동화만을 목표로 하기보다 ``자동 추출 $\rightarrow$ 사용자 확인/수정 $\rightarrow$ 저장''이라는 Human-in-the-loop 흐름을 지원한다.

\newpage

% ============================================================
% 4 평가
% ============================================================
\section{연구 결과 분석 및 평가}

\subsection{평가 목적}

평가는 두 가지 질문에 초점을 둔다.

첫째, 정형 패턴 기반 추출과 LLM 기반 추출은 어떤 종류의 필드에서 서로 다른 성능을 보이는가?

둘째, 두 방식의 처리 시간과 LLM 토큰 사용량에는 어느 정도 차이가 있는가?

본 실험은 프로젝트에 포함된 \texttt{evaluate\_extraction.py}를 사용하였다. 카테고리 분류 자체는 실험 대상에서 제외하고 OCR 텍스트에서 세부 정보를 추출하는 능력에 집중하였다.

\subsection{평가 데이터 및 방법}

평가 스크립트는 \texttt{dataset.csv}에서 OCR 텍스트 길이가 50자 이하인 항목을 제외한 뒤 네 카테고리에서 각각 최대 50개를 random seed 42로 샘플링한다. 따라서 본 실험은 SCHEDULE, PLACE, WISHLIST, MEMO 각각 50개, 총 200개를 대상으로 한다.

비교 대상은 다음과 같다.

\begin{itemize}
    \item \textbf{Regex baseline:} 날짜, 가격, 장소 등에 대한 정규표현식 기반 추출
    \item \textbf{LLM:} GPT-4o-mini를 이용한 카테고리별 정보 추출
    \item \textbf{Reference annotation:} GPT-4o가 동일 OCR 텍스트에서 명시적 정보를 추출하여 생성한 기준값
\end{itemize}

여기서 GPT-4o 기준값은 사람이 독립적으로 작성한 정답지가 아니므로 본 보고서에서는 ``ground truth'' 대신 ``reference annotation''이라고 표현한다.

\subsection{평가 기준}

평가 스크립트는 문자열을 영문/숫자/한글 중심으로 정규화한 뒤 예측값과 기준값 사이에 포함 관계가 있는지를 검사한다. 따라서 정확한 문자열 동일성만 평가하는 것이 아니라 부분 매칭을 허용한다.

이 방식은 OCR 결과의 표현 차이를 어느 정도 흡수할 수 있지만, 동시에 ``문자열 포함''만으로 정답을 판단하므로 의미적으로 완전히 동일한지를 평가하는 방법은 아니다. 따라서 결과는 ``reference annotation 대비 매칭률''로 해석한다.

\subsection{카테고리별 처리 시간}

\begin{table}[H]
\centering
\small
\begin{tabularx}{0.9\linewidth}{l r r r}
\toprule
\textbf{카테고리} & \textbf{Regex 평균} & \textbf{LLM 평균} & \textbf{샘플 수}\\
\midrule
SCHEDULE & 0.199 ms & 1.317 s & 50\\
PLACE & 0.123 ms & 0.889 s & 50\\
WISHLIST & 0.135 ms & 0.784 s & 50\\
MEMO & 0.083 ms & 1.230 s & 50\\
\midrule
전체 평균 & 0.135 ms & 1.06 s & 200\\
\bottomrule
\end{tabularx}
\caption{카테고리별 평균 처리 시간}
\end{table}

측정값은 평가 스크립트가 실제 Regex 함수와 LLM 호출 전후에서 기록한 시간이다. 따라서 Regex의 0.135 ms는 OCR이나 이미지 로딩을 포함한 전체 모바일 파이프라인의 지연 시간이 아니다. 마찬가지로 LLM의 1.06 s도 앱 전체 화면 전환 시간을 의미하지 않는다.

두 방식의 시간 차이는 정형 패턴 탐색과 네트워크를 포함한 LLM 호출의 계산 특성 차이를 보여준다. 그러나 이를 근거로 ``전체 앱이 1ms 이내에 동작한다''고 표현해서는 안 된다.

\subsection{SCHEDULE 추출 성능}

\begin{table}[H]
\centering
\small
\begin{tabularx}{0.92\linewidth}{l r r}
\toprule
\textbf{필드} & \textbf{Regex} & \textbf{LLM}\\
\midrule
날짜 & 66.7\% (30/45) & 33.3\% (15/45)\\
가격 & 71.4\% (10/14) & 85.7\% (12/14)\\
장소 & 9.5\% (4/42) & 66.7\% (28/42)\\
\bottomrule
\end{tabularx}
\caption{SCHEDULE의 reference annotation 대비 매칭률}
\end{table}

SCHEDULE의 날짜에서는 Regex가 높은 매칭률을 보였다. 날짜는 \texttt{2026.06.28}, \texttt{2026년6월24일}과 같이 비교적 규칙적인 문자열 패턴을 갖기 때문이다.

반면 장소에서는 LLM의 매칭률이 크게 높았다. 장소명은 ``장소:'', ``교환처:''와 같은 명시적인 키워드가 존재하지 않을 경우 단순 Regex만으로 추출하기 어렵기 때문이다. 가격 역시 LLM이 더 높은 매칭률을 보였으나, 가격은 Regex도 상당한 수준의 추출이 가능했다.

따라서 SCHEDULE 내부에서도 ``날짜''와 ``장소''는 동일한 알고리즘으로 처리할 필요가 없음을 확인할 수 있다.

\subsection{PLACE 추출 성능}

\begin{table}[H]
\centering
\small
\begin{tabularx}{0.92\linewidth}{l r r}
\toprule
\textbf{필드} & \textbf{Regex} & \textbf{LLM}\\
\midrule
상호명 & 5.6\% (2/36) & 80.6\% (29/36)\\
지역/주소 & -- & 89.7\% (35/39)\\
평점 & -- & 88.2\% (15/17)\\
가격 & 30.8\% (4/13) & 61.5\% (8/13)\\
\bottomrule
\end{tabularx}
\caption{PLACE의 reference annotation 대비 매칭률}
\end{table}

PLACE에서는 Regex의 한계가 더욱 뚜렷하게 나타난다. 평가 스크립트의 Regex 장소 추출은 ``장소'', ``교환처'', ``위치'', ``가게명'', ``호텔'', ``숙소'', ``매장'' 등의 명시적 키워드 주변 문자열이나 \texttt{at}, \texttt{@} 패턴을 대상으로 한다. 따라서 화면에 자연스럽게 나타난 ``성수동 카페 추천''과 같은 텍스트에서 특정 상호명을 의미적으로 식별하는 것은 어렵다.

LLM은 이러한 문맥을 이용할 수 있으므로 상호명과 지역/주소에서 높은 매칭률을 보였다. 이는 LLM을 사용해야 하는 대표적인 정보 유형을 보여준다.

한편 평가 결과에 ``평점''이 포함되어 있지만, 현재 서버 validator의 \texttt{\_enrich\_place\_fields}에서는 \texttt{rating} 필드를 제거한다. 따라서 실험에서 측정된 평점 추출 성능과 최종 앱 데이터 구조에 저장되는 필드는 구분하여 해석해야 한다.

\subsection{WISHLIST 추출 성능}

\begin{table}[H]
\centering
\small
\begin{tabularx}{0.92\linewidth}{l r r}
\toprule
\textbf{필드} & \textbf{Regex} & \textbf{LLM}\\
\midrule
가격 & 87.8\% (43/49) & 95.9\% (47/49)\\
상품명 & -- & 64.3\% (27/42)\\
브랜드 & -- & 38.9\% (7/18)\\
\bottomrule
\end{tabularx}
\caption{WISHLIST의 reference annotation 대비 매칭률}
\end{table}

WISHLIST의 가격은 Regex도 87.8\%의 높은 매칭률을 보였다. 상품 가격은 통화 기호나 ``원''과 같은 패턴이 비교적 명확하기 때문이다.

반면 상품명과 브랜드는 자연어 및 고유명사에 가까우므로 Regex baseline에서 별도의 추출 로직이 없어 비교가 어렵다. LLM은 상품명 64.3\%, 브랜드 38.9\%의 매칭률을 보였으며, 브랜드가 이미지 로고 형태로만 존재하거나 OCR에 포함되지 않는 경우에는 텍스트 기반 LLM만으로 복원하기 어려운 한계도 확인할 수 있다.

\subsection{MEMO 추출 성능}

\begin{table}[H]
\centering
\small
\begin{tabularx}{0.92\linewidth}{l r r}
\toprule
\textbf{필드} & \textbf{Regex} & \textbf{LLM}\\
\midrule
날짜 & 0.0\% (0/2) & 100.0\% (2/2)\\
제목/핵심 정보 & -- & 47.9\% (23/48)\\
\bottomrule
\end{tabularx}
\caption{MEMO의 reference annotation 대비 매칭률}
\end{table}

MEMO의 날짜 평가 대상은 기준값에 날짜가 존재한 2건뿐이므로 100\%라는 수치를 일반적인 날짜 추출 성능으로 확대해서 해석할 수 없다. 표본 수가 작기 때문이다.

제목/핵심 정보는 비정형 텍스트의 요약 및 제목 생성 문제이므로 Regex baseline과 직접 비교하기 어렵다. LLM의 47.9\% 매칭률은 해당 평가 로직에서 reference annotation과 부분적으로 일치한 비율이며, ``메모 요약의 절대 정확도''를 의미하지 않는다.

\subsection{종합 성능 비교}

\begin{table}[H]
\centering
\small
\begin{tabularx}{\linewidth}{l Y Y}
\toprule
\textbf{정보 유형} & \textbf{Regex의 특성} & \textbf{LLM의 특성}\\
\midrule
정형 날짜 & 빠르고 일정한 패턴에 강함 & 문맥 해석에 따른 누락 가능\\
정형 가격 & 높은 추출 가능 & OCR 오류 및 문맥 보정 가능\\
장소/상호명 & 명시적 키워드가 있을 때 제한적 & 문맥 기반 식별에 강함\\
상품명 & 별도 의미 분석 없이는 제한적 & 자연어 문맥에서 추출 가능\\
주소/지역 & 다양한 표현 때문에 제한적 & 구조화 및 정규화와 결합 가능\\
제목/요약 & 일반적으로 부적합 & 의미 기반 처리 가능\\
\bottomrule
\end{tabularx}
\caption{정보 유형별 Regex와 LLM의 특성}
\end{table}

본 결과에서 중요한 것은 ``LLM이 항상 우수하다''는 결론이 아니다. 오히려 필드의 특성에 따라 적합한 처리 방식이 달라진다는 점이다. 날짜나 가격처럼 패턴이 명확한 데이터는 로컬 규칙 기반 처리의 장점이 크고, 장소명이나 상품명처럼 의미 해석이 필요한 데이터는 LLM의 장점이 나타난다.

따라서 솔티의 하이브리드 구조는 단일 모델의 우열을 전제로 하지 않고 서로 다른 처리기의 장점을 결합한다는 관점에서 설명하는 것이 적절하다.

\subsection{LLM 토큰 사용량 및 비용}

평가 데이터에 기록된 LLM 평균 토큰 사용량은 카테고리별로 다음과 같다.

\begin{table}[H]
\centering
\small
\begin{tabularx}{0.85\linewidth}{l r}
\toprule
\textbf{카테고리} & \textbf{평균 토큰/요청}\\
\midrule
SCHEDULE & 452.12\\
PLACE & 354.28\\
WISHLIST & 332.18\\
MEMO & 314.20\\
\midrule
전체 평균 & 약 363\\
\bottomrule
\end{tabularx}
\caption{카테고리별 LLM 토큰 사용량}
\end{table}

평가 문서에는 200건 처리에 총 72,639 tokens, 총 비용 약 \$0.044, 건당 약 0.3원이라는 기록이 있다. 다만 API 가격은 사용 시점과 모델 가격 정책에 따라 달라질 수 있으므로 최종보고서에서는 이 수치를 ``2026년 8월 23일 평가 당시 측정된 비용''으로 명시하고, 향후 운영비를 고정값으로 주장하지 않는다.

\subsection{다중 처리 및 서버 동시성}

클라이언트는 다중 이미지 요청을 최대 20개 단위로 분할한다. 서버의 Semaphore는 동시에 실행되는 LLM 호출을 3개로 제한한다.

\placeholder{그림 삽입: 단건 요청과 다중 요청의 처리 흐름 및 Semaphore=3 구조}

기존 초안에는 20건 처리 시간이 10.77초라는 값이 있었으나, 현재 확인한 \texttt{evaluation\_results\_v2.json}은 단건 Regex/LLM 추출 시간 평가 자료이며 이 값을 직접 검증할 수 있는 동일 실험 결과가 포함되어 있지 않다. 따라서 최종본에서는 해당 수치를 확정적으로 사용하지 않고, 실제 배치 벤치마크를 다시 수행한 후 입력하도록 한다.

\placeholder{실험 필요: 20건 일괄 처리의 전체 wall-clock time, 실패율, 평균/최대 latency를 동일 환경에서 3회 이상 반복 측정}

\subsection{개인정보 보호 평가}

현재 코드에는 다양한 개인정보 탐지 규칙과 NER 모델이 존재하지만, ``50개의 모의 개인정보 스크린샷에서 100\% 차단''과 같은 독립적인 정량 실험 결과는 현재 확인된 평가 JSON에 포함되어 있지 않다.

따라서 본 최종보고서에서는 탐지 로직 자체는 구현 사실로 기술하되, 정량적인 탐지율을 임의로 작성하지 않는다.

\begin{table}[H]
\centering
\small
\begin{tabularx}{\linewidth}{l l l}
\toprule
\textbf{평가 대상} & \textbf{구현 여부} & \textbf{정량 평가}\\
\midrule
주민등록번호 Regex & 구현 & [실험 필요]\\
카드번호 Regex & 구현 & [실험 필요]\\
전화번호 Regex & 구현 & [실험 필요]\\
계좌번호 Regex & 구현 & [실험 필요]\\
여권/MRZ 탐지 & 구현 & [실험 필요]\\
운전면허번호 탐지 & 구현 & [실험 필요]\\
NER PER/PHONE & 구현 & [실험 필요]\\
서버 요청의 ENC 토큰 전달 & 구현 & [페이로드 로그 검증 필요]\\
\bottomrule
\end{tabularx}
\caption{개인정보 보호 기능의 구현 및 평가 현황}
\end{table}

개인정보 보호 성능을 최종적으로 정량화하려면 각 유형별 정상 입력, 공백/하이픈 변형, OCR 오류, 경계 사례를 포함한 별도의 테스트셋을 구축하고 precision, recall, false negative rate를 측정해야 한다.

\subsection{개발 과정의 문제와 개선}

프로젝트 코드를 기준으로 확인되는 주요 개선 사례는 다음과 같다.

\subsubsection{기프티콘 OCR 텍스트 축약 문제}

초기에는 기프티콘 텍스트에서 교환처, 상품명, 유효기간을 추출하여 간략한 세 줄의 텍스트로 만드는 방식이 존재했다. 그러나 이 방식은 원본 OCR의 다른 정보가 제거되어 LLM이 충분한 문맥을 볼 수 없고, 결과적으로 예제에 포함된 상호명을 잘못 사용하는 문제가 발생할 수 있었다.

현재 코드에서는 해당 축약 결과를 LLM 입력으로 사용하지 않고 원본 OCR 전체를 보존하도록 수정되어 있다. 이는 단순 기능 추가가 아니라 ``전처리가 후속 AI의 정보량을 과도하게 줄일 수 있다''는 문제를 해결한 사례이다.

\subsubsection{지역명 정규화 문제}

장소 정보에서는 지역 필터링을 위해 표준 지역명을 사용한다. 특히 ``해운대''라는 문자열이 ``대구''와 부분적으로 겹칠 수 있는 문제를 별도의 예외로 처리하고, 경상남도/경상북도/전라남도 등의 표현을 경남/경북/전라 등으로 정규화한다.

이러한 규칙은 LLM 출력 자체를 신뢰하는 것이 아니라, 애플리케이션이 실제로 필요로 하는 데이터 규격에 맞추는 후처리 계층의 역할을 보여준다.

\subsubsection{LLM 출력의 불확실성}

LLM 응답은 생성 결과이므로 필드 누락이나 허용되지 않은 타입이 반환될 가능성이 있다. validator는 이를 애플리케이션 수준에서 검사하고 누락 필드를 표시한다.

따라서 본 시스템은 ``LLM의 응답을 곧바로 DB에 저장''하지 않고, 검증 및 enrichment를 거친 결과를 사용하는 구조로 개선되었다.

\newpage

% ============================================================
% 5 결론
% ============================================================
\section{결론 및 향후 연구 방향}

\subsection{연구 결과 요약}

본 연구에서는 스마트폰 스크린샷을 일정, 장소, 위시리스트, 메모 등 활용 가능한 구조화 데이터로 변환하는 솔티 애플리케이션을 구현하였다.

시스템은 Flutter 기반 모바일 클라이언트에서 Google ML Kit OCR을 수행하고, TFLite 기반 온디바이스 텍스트 분류기로 네 가지 대분류를 판단한다. 분류기의 입력은 프로젝트의 vocabulary와 IDF를 이용한 TF-IDF character n-gram 특징으로 구성된다.

개인정보 보호 측면에서는 NER와 Regex를 결합하여 전화번호, 주민등록번호, 카드번호, 계좌번호, 예약번호, 여권 및 MRZ 등 다양한 패턴을 탐지하고, 탐지된 문자열을 AES-256-CBC로 암호화하여 \texttt{[ENC:...]} 토큰으로 변환한다. 이를 통해 서버의 LLM이 원문의 민감 문자열을 직접 입력받지 않는 구조를 구현하였다.

서버에서는 FastAPI가 분석 요청을 처리하고, 카테고리별 LLM 프롬프트를 이용하여 일정, 장소, 상품, 메모 등의 필드를 구조화한다. Validator는 필수 필드와 타입을 확인하고 지역명, 가격, 알림 설정, 이미지 보존 정책 등 애플리케이션에 필요한 후처리를 수행한다.

저장 계층에서는 SQLite와 JSON을 결합하여 카테고리별로 서로 다른 필드를 유연하게 저장한다. 또한 검색 및 필터링, 이미지 해시 기반 결과 조회, 알림 스케줄러, iCalendar 생성, Kakao Calendar 연동 등의 후속 기능을 구현하였다.

\subsection{정량 평가의 의의}

200개 샘플을 이용한 정보 추출 평가에서 전체 평균 Regex 측정 시간은 0.135 ms, LLM 측정 시간은 약 1.06 s였다. Regex는 날짜와 가격처럼 정형적인 패턴에서 높은 매칭률을 보였으며, LLM은 장소명, 상호명, 주소 등 문맥 기반 정보에서 더 높은 매칭률을 보였다.

특히 SCHEDULE의 장소 매칭은 Regex 9.5\% 대비 LLM 66.7\%, PLACE의 상호명은 Regex 5.6\% 대비 LLM 80.6\%로 나타났다. 이러한 결과는 정형 정보와 비정형 정보에 동일한 추출기를 적용하는 것보다 정보 특성에 따라 역할을 분리하는 것이 합리적임을 보여준다.

다만 본 평가의 reference annotation은 GPT-4o가 생성한 기준값이며 인간 전문가가 독립적으로 작성한 정답지가 아니다. 따라서 결과를 절대적인 정확도로 표현하기보다 ``GPT-4o reference annotation 대비 매칭률''로 해석해야 한다.

\subsection{연구의 한계}

첫째, OCR 자체의 정확도를 독립적으로 평가한 실험은 본 정보 추출 평가와 별개의 문제이다. OCR 오류가 발생하면 이후 분류와 구조화 단계에도 영향을 줄 수 있다.

둘째, 현재 평가셋은 카테고리별 50건, 총 200건으로 제한되어 있다. 다양한 앱 UI, 해상도, 언어 혼합, OCR 오류 유형을 충분히 대표한다고 보기 어렵다.

셋째, 개인정보 보호 로직은 코드상 다양한 패턴을 지원하지만 각 유형의 precision과 recall을 정량적으로 검증한 독립적인 보안 평가셋이 부족하다.

넷째, AES-256-CBC를 사용하고 있어 암호화된 데이터의 무결성 검증을 별도로 제공하지 않는다. 따라서 향후 authenticated encryption 방식으로 개선할 필요가 있다.

다섯째, 서버 기반 LLM을 사용하므로 네트워크 연결이 필요한 후처리 단계가 존재한다. 온디바이스 분류와 마스킹은 로컬에서 수행되지만 상세 구조화까지 완전한 오프라인으로 동작하는 것은 아니다.

여섯째, 로컬 DB 저장 경로 중 일부에는 암호화된 상태가 아닌 구조화 텍스트를 저장하는 코드가 존재한다. 따라서 ``로컬 DB의 모든 개인정보가 항상 암호화되어 저장된다''고 주장해서는 안 된다.

\subsection{향후 연구 방향}

\subsubsection{개인정보 보호 평가의 정량화}

향후에는 개인정보 유형별 테스트셋을 구축하고 NER, Regex, Hybrid 방식의 precision, recall, F1-score를 비교할 필요가 있다. 특히 OCR에서 하이픈이나 공백이 사라진 경우, 숫자가 일부 오인식된 경우, 민감정보와 일반 숫자가 함께 존재하는 경우 등을 포함해야 한다.

\subsubsection{인증 암호화 적용}

현재 AES-CBC 기반 암호화는 기밀성 중심의 구현이다. 향후 AES-GCM과 같은 AEAD 방식으로 전환하여 암호문에 대한 무결성 검증까지 제공할 수 있다.

\subsubsection{온디바이스 구조화 모델}

현재 시스템은 상세 구조화를 클라우드 LLM에 위임한다. 향후 모바일 환경에서 충분한 성능을 갖는 소형 언어 모델을 적용하면 네트워크가 없는 환경에서도 전체 파이프라인을 실행할 수 있다. 다만 모델 크기, 배터리, 메모리, 발열 및 응답 속도에 대한 별도의 평가가 필요하다.

\subsubsection{더 큰 평가셋과 인간 평가}

GPT-4o reference annotation을 인간 정답으로 대체하거나, 최소한 일부 표본에 대해 사람의 독립적인 검수를 수행하면 보다 신뢰도 높은 정확도 평가가 가능하다. 또한 200건보다 큰 데이터셋에서 카테고리와 UI 유형별 성능을 비교할 필요가 있다.

\subsubsection{End-to-end latency 측정}

현재 평가의 0.135 ms는 Regex 단계만을 측정한다. 향후에는 이미지 선택부터 OCR, 전처리, 분류, 개인정보 보호, 네트워크 요청, LLM 응답, validator, DB 저장까지 포함한 end-to-end latency를 측정해야 한다.

\subsubsection{스토리지 정책 검증}

\texttt{keep\_photo} 정책과 실제 파일 삭제/보존 동작을 분리하여 테스트하고, 사용자가 원본 보존 여부를 직접 선택할 수 있는 기능을 추가하는 것도 고려할 수 있다.

\newpage

% ============================================================
% 6 참고문헌
% ============================================================
\section{참고 문헌}

\begin{enumerate}[label={[\arabic*]}]
    \item 이종한, 박성재, 이채은, 권아현, 김민재, 신현석, 황연우, 서왕덕,
    ``개인정보 보호형 LLM 기반 실시간 근로계약서 분석 및 가이드 시스템,''
    \textit{2026 한국정보기술학회 하계 종합학술대회 논문집}, 2026.

    \item 정휘수, 이지훈, 주성진, 이길호,
    ``On-device LLM모델의 반응성 향상을 위한 동적 배치 최적화 기술,''
    \textit{2024 한국소프트웨어종합학술대회 논문집}, 2024.

    \item B. Chen, Y. Chen, X. Jiang, et al.,
    ``Unleashing the Potential of Prompt Engineering in Large Language Models: A Comprehensive Review,''
    arXiv preprint arXiv:2310.14735, 2023.

    \item Google,
    ``ML Kit Text Recognition Documentation,''
    Google Developers.
    \url{https://developers.google.com/ml-kit/vision/text-recognition}

    \item Flutter,
    ``Flutter Documentation,''
    \url{https://docs.flutter.dev/}

    \item Pydantic,
    ``Pydantic Documentation,''
    \url{https://docs.pydantic.dev/}

    \item OpenAI,
    ``OpenAI API Documentation,''
    \url{https://platform.openai.com/docs/}

    \item SQLite,
    ``SQLite Documentation,''
    \url{https://www.sqlite.org/docs.html}
\end{enumerate}

\newpage

% ============================================================
% 부록 A
% ============================================================
\section*{부록 A. 주요 소스 파일과 구현 기능}
\addcontentsline{toc}{section}{부록 A. 주요 소스 파일과 구현 기능}

\begingroup
\renewcommand{\arraystretch}{1.25}
\begin{longtable}{>{\raggedright\arraybackslash\small}p{0.44\linewidth} >{\raggedright\arraybackslash\small}p{0.51\linewidth}}
\toprule
\textbf{소스 파일} & \textbf{주요 기능}\\
\midrule
\texttt{lib/main.dart} & Flutter 앱 초기화, OCR/분류/분석 흐름, UI 및 저장 연계\\
\texttt{lib/\allowbreak core/\allowbreak ml/\allowbreak on\_device\_\allowbreak text\_classifier.dart} & TF-IDF character n-gram 전처리 및 TFLite 분류\\
\texttt{lib/\allowbreak core/\allowbreak utils/\allowbreak masking\_helper.dart} & 개인정보 탐지 및 암호화 토큰 치환\\
\texttt{lib/\allowbreak core/\allowbreak utils/\allowbreak ner\_classifier.dart} & TFLite NER 모델 추론\\
\texttt{lib/\allowbreak core/\allowbreak utils/\allowbreak ner\_tokenizer.dart} & WordPiece 토큰화 및 offset 관리\\
\texttt{lib/\allowbreak core/\allowbreak utils/\allowbreak crypto\_helper.dart} & AES-256-CBC 암호화/복호화\\
\texttt{lib/\allowbreak core/\allowbreak storage/\allowbreak database\_helper.dart} & SQLite CRUD\\
\texttt{lib/\allowbreak core/\allowbreak storage/\allowbreak app\_storage.dart} & 카테고리별 앱 저장 디렉터리 초기화\\
\texttt{lib/\allowbreak services/\allowbreak kakao\_calendar\_\allowbreak service.dart} & Kakao Calendar 인증 및 일정 생성\\
\texttt{app/\allowbreak routers/\allowbreak analyze.py} & 분석 API 라우팅\\
\texttt{app/\allowbreak llm\_client.py} & LLM 호출 및 실패 처리\\
\texttt{app/\allowbreak prompts.py} & 분류 및 카테고리별 구조화 프롬프트\\
\texttt{app/\allowbreak validator.py} & 결과 검증 및 enrichment\\
\texttt{app/\allowbreak concurrency.py} & LLM 동시성 제한\\
\texttt{app/\allowbreak database.py} & 서버 SQLite 저장/검색/알림 데이터 관리\\
\texttt{app/\allowbreak calendar.py} & iCalendar 문자열 생성\\
\texttt{app/\allowbreak scheduler.py} & 일정/만료 알림 생성\\
\texttt{evaluate\_\allowbreak extraction.py} & Regex와 LLM 추출 성능 평가\\
\texttt{evaluation\_\allowbreak results\_\allowbreak v2.json} & 200개 샘플의 평가 결과\\
\bottomrule
\end{longtable}
\endgroup

\newpage

% ============================================================
% 부록 B
% ============================================================
\section*{부록 B. 최종보고서 작성 시 추가 실험 항목}
\addcontentsline{toc}{section}{부록 B. 최종보고서 작성 시 추가 실험 항목}

\begin{enumerate}
    \item \textbf{End-to-end latency}
    \begin{itemize}
        \item 이미지 선택
        \item OCR
        \item 온디바이스 분류
        \item 개인정보 탐지 및 암호화
        \item 네트워크 전송
        \item LLM
        \item validator
        \item DB 저장
    \end{itemize}

    \item \textbf{개인정보 탐지 성능}
    \begin{itemize}
        \item 전화번호
        \item 카드번호
        \item 계좌번호
        \item 주민등록번호
        \item 여권번호/MRZ
        \item 운전면허번호
        \item 예약번호
        \item 쿠폰/바코드 번호
    \end{itemize}

    \item \textbf{다중 이미지 성능}
    \begin{itemize}
        \item 1건 / 5건 / 10건 / 20건
        \item 평균 wall-clock time
        \item p95 latency
        \item 실패율
        \item 재시도 횟수
    \end{itemize}

    \item \textbf{분류 모델 성능}
    \begin{itemize}
        \item Confusion Matrix
        \item Precision / Recall / F1
        \item 클래스별 샘플 수
    \end{itemize}

    \item \textbf{사용자 평가}
    \begin{itemize}
        \item 자동 추출 결과 수정 횟수
        \item 검색 성공률
        \item 일정 등록 성공률
        \item 사용 전후 처리 시간
    \end{itemize}
\end{enumerate}

\placeholder{그림 삽입: 최종 시스템 UI 화면 3--5장}

\placeholder{표 삽입: 실제 사용자 시나리오별 처리 결과}

\end{document}
